\documentclass[10pt]{article}

% ---------- IDIOMA Y CODIFICACIÓN ----------
\usepackage[utf8]{inputenc}
\usepackage[spanish]{babel}
\usepackage{amsmath}
\usepackage[margin=2.5cm]{geometry}
\usepackage{setspace}
\usepackage{graphicx}
\usepackage{listings}
\usepackage{minted}
\usepackage{tikz}
\usetikzlibrary{circuits.logic.US,babel}

% ---------- TABLAS (Actividad 1: identificación de dispositivos) ----------
\usepackage{booktabs}   % líneas de tabla profesionales (\toprule, \midrule, \bottomrule)
\usepackage{array}
\usepackage{ragged2e}
\usepackage[table]{xcolor} % opcional, para color en encabezado
\usepackage{tabularx}   % columnas que se ajustan al ancho de página
\usepackage{longtable}  % tablas que pueden partirse entre páginas (útil si tu tabla de dispositivos es larga)
\usepackage{caption}    % control de estilo de "Tabla X: ..."
\usepackage{float}
\usepackage{makecell}

% ---------- GLOSARIO (Actividad 2) ----------
% Enfoque simple con "description", sin necesidad de correr makeglossaries aparte
% (evita agregar un paso extra a tu pipeline de compilación).
\usepackage{enumitem}
\newlist{glosario}{description}{1}
\setlist[glosario]{
    font=\normalfont\bfseries,
    leftmargin=0pt,
    labelsep=0.5em,
    itemsep=0.5em
}

% ---------- Definiciones auxiliares ----------
\lstdefinestyle{asm}{
  language={[x86masm]Assembler},
  basicstyle=\ttfamily\small,
  keywordstyle=\bfseries,
  commentstyle=\itshape,
  numbers=left,
  numberstyle=\tiny,
  numbersep=6pt,
  frame=single,
  breaklines=true,
  showstringspaces=false,
  tabsize=4
}
%---- MINTED..----- 

\usepackage{xcolor}
\usepackage{minted}
\definecolor{bg}{rgb}{0.90,0.91,0.83}

\providecommand{\espacioimagen}[1]{%
  \fbox{\parbox[c][4.5cm][c]{0.9\linewidth}{\centering\itshape Insertar captura: #1}}}
 
% \figuraimagen{archivo.png}{Leyenda}{etiqueta}{descripción del recuadro}
\providecommand{\figuraimagen}[4]{%
  \begin{figure}[htbp]
    \centering
    \IfFileExists{img/#1}{\includegraphics[width=0.9\linewidth]{img/#1}}{\espacioimagen{#4\\[2pt]{\upshape\small\texttt{img/#1}}}}
    \caption{#2}\label{#3}
  \end{figure}}
 

%      \relleno{Tu texto}    -> muestra tu texto en rojo (útil para ir dejando borradores)
\newcommand{\relleno}[1][\textbf{[COMPLETAR AQUÍ]}]{\textcolor{red}{#1}}
% ---------- CONFIGURACIONES PARA ENCABEZADO Y PORTADA ----------
\usepackage{colortbl}
\usepackage[colorlinks=true,linkcolor=black,urlcolor=black,citecolor=black]{hyperref}

% Definición de colores y columnas personalizadas (usadas en la portada)
\definecolor{tablebackground}{rgb}{0.8, 0, 0}
\newcolumntype{x}[1]{>{\centering\arraybackslash\hspace{0pt}}p{#1}}
\newcolumntype{M}[1]{>{\centering\arraybackslash}m{#1}}
\newcolumntype{N}{@{}m{0pt}@{}}
\newcolumntype{P}[1]{>{\RaggedRight\arraybackslash}p{#1}}

% Variables de datos para la portada y encabezado
\newcommand{\itemEmail}{ccapatinta@unsa.edu.pe}
\newcommand{\itemStudent}{Camilo B.A. Capatinta Almanza (Teoría grupo B)}
\newcommand{\itemCourse}{Lab. Arq. de computadoras}
\newcommand{\itemCourseCode}{1702227}
\newcommand{\itemSemester}{IV}
\newcommand{\itemUniversity}{Universidad Nacional de San Agustín de Arequipa}
\newcommand{\itemFaculty}{Facultad de Ingeniería de Producción y Servicios}
\newcommand{\itemDepartment}{Departamento Académico de Ingeniería de Sistemas e Informática}
\newcommand{\itemSchool}{Escuela Profesional de Ingeniería de Sistemas}
\newcommand{\itemAcademic}{2026 - B}
\newcommand{\itemInput}{20/10/2026}
\newcommand{\itemOutput}{--/10/2026}
\newcommand{\itemPracticeNumber}{05}
\newcommand{\itemTheme}{Programación para el microprocesador INTEL - Instrucciones Aritméticas}

% Configuración del encabezado y pie de página (fancyhdr)
\usepackage{fancyhdr}
\pagestyle{fancy}
\fancyhf{}
\setlength{\headheight}{40.51407pt}
\renewcommand{\headrulewidth}{1pt}
\renewcommand{\footrulewidth}{1pt}

% Encabezado
\fancyhead[L]{\raisebox{-0.2\height}{\includegraphics[width=3cm]{img/logo_episunsa.png}}}
\fancyhead[C]{\fontsize{7}{7}\selectfont    \itemUniversity \\ \itemFaculty \\ \itemDepartment \\ \itemSchool \\ \textbf{\itemCourse}}
\fancyhead[R]{\raisebox{-0.2\height}{\includegraphics[width=1.2cm]{img/logo_abet.png}}}

% Pie de página
\fancyfoot[L]{Estudiante Camilo Capatinta}
\fancyfoot[C]{\itemCourse}
\fancyfoot[R]{Página \thepage}

\begin{document}

% --- INICIO DE LA PORTADA ---
	\vspace*{10px}
	
	\begin{center}
		\parbox{0.8\textwidth}{\centering\fontsize{17}{20}\selectfont\textbf{Informe de laboratorio \itemPracticeNumber}}
	\end{center}
	\begin{center}
		\parbox{0.8\textwidth}{\centering\textbf{\Large Tema: \itemTheme}}
	\end{center}

\begin{flushright}
    \begin{tabular}{|M{2.5cm}|N|}
        \hline 
        \rowcolor{tablebackground}
        \color{white} \textbf{Nota}  \\
        \hline 
             \\[30pt]
        \hline          
    \end{tabular}
\end{flushright}    

\begin{table}[h!]
    \begin{tabular}{|x{4.7cm}|x{4.8cm}|x{4.8cm}|}        
        \hline 
        \rowcolor{tablebackground}
        \color{white} \textbf{Estudiante} & \color{white}\textbf{Escuela}  & \color{white}\textbf{Asignatura}   \\
        \hline 
        {\itemStudent \par \itemEmail} & \itemSchool & {\itemCourse \par Semestre: \itemSemester \par Código: \itemCourseCode}     \\
        \hline          
    \end{tabular}
\end{table}     

\begin{table}[H]
    \begin{tabular}{|x{4.7cm}|x{4.8cm}|x{4.8cm}|}
        \hline 
        \rowcolor{tablebackground}
        \color{white}\textbf{Laboratorio} & \color{white}\textbf{Tema}  & \color{white}\textbf{Duración}   \\
        \hline 
        \itemPracticeNumber & \itemTheme & 02 horas   \\
        \hline 
    \end{tabular}
\end{table}

\begin{table}[H]
    \begin{tabular}{|x{4.7cm}|x{4.8cm}|x{4.8cm}|}
        \hline 
        \rowcolor{tablebackground}
        \color{white}\textbf{Semestre académico} & \color{white}\textbf{Fecha de inicio}  & \color{white}\textbf{Fecha de entrega}   \\
        \hline 
        \itemAcademic & \itemInput &  \itemOutput  \\
        \hline 
    \end{tabular}
\end{table}

\newpage
\tableofcontents
\thispagestyle{plain}
\clearpage

\section{Objetivos}
 
\textbf{Objetivo general.} Analizar la forma en que se ejecutan algunas instrucciones aritméticas y de manejo de la pila.

\textbf{Objetivos específicos.}

\begin{itemize}
  \item Observar el uso de algunas instrucciones aritméticas y las banderas que son afectadas.
  \item Identificar las banderas que afecta cada instrucción observada.
  \item Utilizar correctamente los registros del microprocesador al realizar operaciones aritméticas, tanto de 8 bits como de 16 bits.
  \item Analizar detenidamente el proceso de cada operación.
\end{itemize}
 
% ============================================================
\section{Materiales y equipos}
 
\begin{itemize}
  \item Guía de laboratorio.
  \item Computadora PC con el programa EMU8086.
\end{itemize}
 
% ============================================================
\section{Actividades}
 
\subsection{Preparación del entorno}
 
\begin{itemize}
  \item Se ejecutó el emulador EMU8086 haciendo doble clic en su icono del escritorio y se creó un nuevo documento tipo COM.
  \item Cada programa se escribió con su encabezado \texttt{\#make\_COM\#} y \texttt{org 100h}, necesario para generar un archivo \texttt{.COM}.
  \item Para observar el efecto de cada instrucción se usó la ejecución paso a paso (\textit{single step}), junto con las ventanas \textbf{flags} y \textbf{memory} del emulador.
\end{itemize}
 
Las banderas que se monitorearon se resumen en la Tabla~\ref{tab:banderas} y las instrucciones aritméticas estudiadas en la Tabla~\ref{tab:instrucciones}.
 
\begin{table}[htbp]
  \centering
  \caption{Banderas del registro de estados utilizadas en la práctica.}
  \label{tab:banderas}
  \begin{tabular}{|c|p{3cm}|p{6.5cm}|}
    \hline
    \textbf{Bandera} & \textbf{Nombre} & \textbf{Se activa (=1) cuando\ldots} \\
    \hline
    O & Desbordamiento (\textit{overflow}) & el resultado, interpretado con signo, no cabe en el tamaño del operando. \\
    \hline
    S & Signo & el bit más significativo del resultado es 1 (bit 7 en operaciones de 8 bits, bit 15 en las de 16 bits). \\
    \hline
    Z & Cero & el resultado de la operación es 0. \\
    \hline
    A & Acarreo auxiliar & hay acarreo (suma) o préstamo (resta) entre el bit 3 y el bit 4. \\
    \hline
    P & Paridad & el byte bajo del resultado tiene un número par de unos. \\
    \hline
    C & Acarreo (\textit{carry}) & hay acarreo fuera del bit más significativo (suma) o préstamo (resta), interpretando los datos sin signo. \\
    \hline
  \end{tabular}
\end{table}
 
\begin{table}[htbp]
  \centering
  \caption{Instrucciones aritméticas estudiadas.}
  \label{tab:instrucciones}
  \begin{tabular}{|l|p{4.2cm}|p{5.6cm}|}
    \hline
    \textbf{Instrucción} & \textbf{Operación} & \textbf{Banderas y observaciones} \\
    \hline
    \texttt{ADD} & destino $\leftarrow$ destino $+$ fuente & Modifica O, S, Z, A, P y C. Ignora el acarreo previo. \\
    \hline
    \texttt{ADC} & destino $\leftarrow$ destino $+$ fuente $+$ C & Modifica O, S, Z, A, P y C. Suma el acarreo de la operación anterior. \\
    \hline
    \texttt{SUB} & destino $\leftarrow$ destino $-$ fuente & Modifica O, S, Z, A, P y C. C = 1 indica préstamo. \\
    \hline
    \texttt{SBB} & destino $\leftarrow$ destino $-$ fuente $-$ C & Modifica O, S, Z, A, P y C. Resta el préstamo de la operación anterior. \\
    \hline
    \texttt{MUL} & AX $\leftarrow$ AL $\times$ fuente (8 bits, sin signo) & Solo C y O (= 1 si AH $\neq$ 0). S, Z, A y P quedan indefinidas. \\
    \hline
    \texttt{IMUL} & AX $\leftarrow$ AL $\times$ fuente (8 bits, con signo) & Solo C y O (= 1 si el producto no cabe en AL con signo). Las demás quedan indefinidas. \\
    \hline
    \texttt{DIV} & AL $\leftarrow$ AX $\div$ fuente, AH $\leftarrow$ residuo (sin signo) & Todas las banderas quedan indefinidas. Error si el divisor es 0 o el cociente supera FFh. \\
    \hline
    \texttt{IDIV} & AL $\leftarrow$ AX $\div$ fuente, AH $\leftarrow$ residuo (con signo) & Todas las banderas quedan indefinidas. Error si el divisor es 0 o el cociente sale de $-128$ a $127$. \\
    \hline
  \end{tabular}
\end{table}
 
% ------------------------------------------------------------
\subsection{Primer programa: suma con ADD y ADC}
 
\textbf{Objetivo.} Observar la suma de dos números de 1 byte y de dos números de 2 bytes, cómo se actualizan las banderas y qué diferencia hay entre \texttt{ADD} y \texttt{ADC}. Los datos se declaran con la directiva \texttt{EQU} (Tabla~\ref{tab:p1-datos}).
 
\begin{table}[htbp]
  \centering
  \caption{Datos declarados en el primer programa.}
  \label{tab:p1-datos}
  \begin{tabular}{|c|c|c|c|}
    \hline
    \textbf{Constante} & \textbf{Decimal} & \textbf{Hexadecimal} & \textbf{Tamaño} \\
    \hline
    \texttt{byte1} & 125 & 7Dh & 8 bits \\
    \hline
    \texttt{byte2} & 150 & 96h & 8 bits \\
    \hline
    \texttt{word1} & 8355 & 20A3h & 16 bits \\
    \hline
    \texttt{word2} & 22840 & 5938h & 16 bits \\
    \hline
  \end{tabular}
\end{table}
 
\textbf{Funcionamiento.}
\begin{itemize}
  \item \texttt{AL = byte1 + byte2} (8 bits): 125 + 150 = 275, que no cabe en un byte. AL queda en 13h y el noveno bit pasa a la bandera C.
  \item \texttt{AX = word1 + word2} (16 bits): 20A3h + 5938h = 79DBh.
  \item Los resultados se copian a memoria: AL en la dirección 011Eh y AX en 011Fh y 0120h (\textit{little endian}: DBh en 011Fh y 79h en 0120h).
\end{itemize}
 
\begin{figure}[H]
    \centering
    \includegraphics[width=0.86\textwidth]{img/12.png}
    \caption{Código del primer programa con ADD y ADC.}
    \label{fig:primerProgramaLab05}
\end{figure}

\textbf{Procedimiento.}
\begin{itemize}
  \item Se digitó y emuló el programa. Antes de ejecutarlo se anotó el contenido de las direcciones 011Eh, 011Fh y 0120h y el estado de las banderas (Figura~\ref{fig:p1-inicio}).
  \item Se ejecutó paso a paso, anotando en cada momento la memoria y las banderas (Tabla~\ref{tab:p1-tabla1}).
\end{itemize}
 
\begin{figure}[H]
    \centering
    \includegraphics[width=0.86\textwidth]{img/12(1).png}
    \caption{Estado inicial de la memoria (011Eh a 0120h) y de las banderas.}
    \label{fig:primerProgramaLab05}
\end{figure}

\begin{table}[htbp]
  \centering
  \small
  \setlength{\tabcolsep}{4pt}
  \caption{Tabla 1: memoria y banderas durante la ejecución del primer programa (con \texttt{ADD}).}
  \label{tab:p1-tabla1}
  \begin{tabular}{|l|c|c|c|c|c|c|c|c|c|}
    \hline
    \textbf{Momento} & \textbf{011Eh} & \textbf{011Fh} & \textbf{0120h} & \textbf{O} & \textbf{S} & \textbf{Z} & \textbf{A} & \textbf{P} & \textbf{C} \\
    \hline
    Inicio & 00h & 00h & 00h & 0 & 0 & 0 & 0 & 0 & 0 \\
    \hline
    Después de la 1.ª suma & 00h & 00h & 00h & 0 & 0 & 0 & 1 & 0 & 1 \\
    \hline
    Antes de la 2.ª suma & 13h & 00h & 00h & 0 & 0 & 0 & 1 & 0 & 1 \\
    \hline
    Después de la 2.ª suma & 13h & 00h & 00h & 0 & 0 & 0 & 0 & 1 & 0 \\
    \hline
    Después del último \texttt{mov} & 13h & DBh & 79h & 0 & 0 & 0 & 0 & 1 & 0 \\
    \hline
  \end{tabular}
 
  \smallskip
  {\small <<Antes de la 2.ª suma>> es después de ejecutar \texttt{mov [011Eh], al} y \texttt{mov ax, word1}; \texttt{MOV} no modifica banderas.}
\end{table}
 
\begin{figure}[H]
    \centering
    \includegraphics[width=0.86\textwidth]{img/12(2).png}
    \caption{Registros, memoria y banderas después de la primera suma.}
    \label{fig:primerProgramaLab05}
\end{figure}
 
\begin{figure}[H]
    \centering
    \includegraphics[width=0.86\textwidth]{img/12(3).png}
    \caption{Registros, memoria y banderas después de la segunda suma y del último \texttt{mov}.}
    \label{fig:primerProgramaLab05}
\end{figure}

\textbf{Respuestas a las preguntas de la guía.}
\begin{itemize}
  \item \textit{¿Con qué datos se realizan las operaciones y cuáles son los resultados?} Con 7Dh + 96h (125 + 150) en 8 bits, cuyo resultado es 13h con acarreo, y con 20A3h + 5938h en 16 bits, cuyo resultado es 79DBh.
  \item \textit{¿Dónde se almacena el resultado al final?} En memoria: el byte 13h en 011Eh y la palabra 79DBh en 011Fh (DBh) y 0120h (79h). AX conserva además 79DBh.
  \item \textit{¿Concuerda el resultado con lo esperado?} Sí. El resultado de 8 bits es 275 $-$ 256 = 19 = 13h, y el 256 que sobra es el acarreo que se ve en C.
\end{itemize}
 
\subsubsection{Acarreo invertido con \texttt{ADD}}
 
Se volvió a cargar el programa y se invirtió manualmente la bandera C (de 0 a 1) antes de ejecutar (equivale a insertar \texttt{stc} al inicio). Se ejecutó hasta la primera suma y luego hasta la segunda.
 
\begin{table}[htbp]
  \centering
  \small
  \setlength{\tabcolsep}{4pt}
  \caption{Primer programa con \texttt{ADD} y C = 1 al inicio.}
  \label{tab:p1-add-cf}
  \begin{tabular}{|l|c|c|c|c|c|c|c|c|}
    \hline
    \textbf{Instrucción} & \textbf{Operandos} & \textbf{Resultado} & \textbf{O} & \textbf{S} & \textbf{Z} & \textbf{A} & \textbf{P} & \textbf{C} \\
    \hline
    \texttt{add al, byte2} & 7Dh, 96h & 13h & 0 & 0 & 0 & 1 & 0 & 1 \\
    \hline
    \texttt{add ax, word2} & 20A3h, 5938h & 79DBh & 0 & 0 & 0 & 0 & 1 & 0 \\
    \hline
  \end{tabular}
\end{table}
 
\begin{figure}[H]
    \centering
    \includegraphics[width=0.86\textwidth]{img/13.png}
    \caption{Resultados de \texttt{ADD} con el acarreo activado manualmente al inicio.}
    \label{fig: primerP_stc.asm}
\end{figure}
 
\textbf{Observaciones.} El resultado de ambas sumas \textbf{no cambió}: \texttt{ADD} no usa el valor previo de C, solo lo actualiza con el acarreo que produce la propia operación. En la segunda suma C pasó a 0 aunque antes valía 1.
 
\subsubsection{Variante con \texttt{ADC}}
 
Se sustituyeron las dos instrucciones \texttt{ADD} por \texttt{ADC} y se repitió el procedimiento desde la inversión del acarreo:
 
\begin{figure}[H]
    \centering
    \includegraphics[width=0.86\textwidth]{img/13(1).png}
    \caption{Primer programa con \textit{ADC} y \textit{C} = 1 al inicio.}
    \label{fig: primerP_adc.asm}
\end{figure}
 
\begin{table}[htbp]
  \centering
  \small
  \setlength{\tabcolsep}{4pt}
  \caption{Primer programa con \texttt{ADC} y C = 1 al inicio.}
  \label{tab:p1-adc}
  \begin{tabular}{|l|c|c|c|c|c|c|c|c|}
    \hline
    \textbf{Instrucción} & \textbf{Operandos} & \textbf{Resultado} & \textbf{O} & \textbf{S} & \textbf{Z} & \textbf{A} & \textbf{P} & \textbf{C} \\
    \hline
    \texttt{adc al, byte2} & 7Dh + 96h + 1 & 14h & 0 & 0 & 0 & 1 & 1 & 1 \\
    \hline
    \texttt{adc ax, word2} & 20A3h + 5938h + 1 & 79DCh & 0 & 0 & 0 & 0 & 0 & 0 \\
    \hline
  \end{tabular}
\end{table}
 
\begin{figure}[H]
    \centering
    \includegraphics[width=0.86\textwidth]{img/13(2).png}
    \caption{Resultados del primer programa con \textit{ADC}.}
    \label{fig: primerP_adc.asm}
\end{figure}
 
\textbf{Observaciones.} Con \texttt{ADC} ambos resultados aumentaron en 1 respecto a \texttt{ADD}: 13h pasó a 14h y 79DBh pasó a 79DCh, porque la instrucción suma además el acarreo. Las banderas A y P cambiaron porque dependen del resultado: 14h tiene dos unos (P = 1) y 79DCh tiene cinco (P = 0). Aunque se empiece con C = 0, la segunda suma también da 79DCh, ya que recibe el acarreo C = 1 que dejó la primera.
 
\subsubsection{Datos modificados}
 
Se restauró \texttt{ADD} y se cambiaron los datos por: \texttt{byte1 equ 9Dh}, \texttt{byte2 equ 83h}, \texttt{word1 equ 5938} y \texttt{word2 equ -5938} (5938 = 1732h y $-5938$ = E8CEh en complemento a 2).
 
\begin{table}[htbp]
  \centering
  \small
  \setlength{\tabcolsep}{4pt}
  \caption{Primer programa con los datos modificados (\texttt{ADD}).}
  \label{tab:p1-datos-nuevos}
  \begin{tabular}{|l|c|c|c|c|c|c|c|c|}
    \hline
    \textbf{Instrucción} & \textbf{Operandos} & \textbf{Resultado} & \textbf{O} & \textbf{S} & \textbf{Z} & \textbf{A} & \textbf{P} & \textbf{C} \\
    \hline
    \texttt{add al, byte2} & 9Dh, 83h & 20h (120h) & 1 & 0 & 0 & 1 & 0 & 1 \\
    \hline
    \texttt{add ax, word2} & 1732h, E8CEh & 0000h (10000h) & 0 & 0 & 1 & 1 & 1 & 1 \\
    \hline
  \end{tabular}
 
  \smallskip
  {\small Entre paréntesis, el resultado completo incluyendo el acarreo.}
\end{table}
 
\begin{figure}[H]
    \centering
    \includegraphics[width=0.86\textwidth]{img/14.png}
    \caption{Programa inicial con los datos modificados.}
    \label{fig: primerP_modificado.asm}
\end{figure}

\begin{figure}[H]
    \centering
    \includegraphics[width=0.86\textwidth]{img/14(1).png}
    \caption{Registros, memoria y banderas con los datos modificados (emulador con los datos 9Dh, 83h, 5938 y -5938).}
    \label{fig: primerP_modificado.asm}
\end{figure}
 
\textbf{Observaciones.} En memoria quedaron 20h en 011Eh y 0000h en 011Fh y 0120h. Las banderas reaccionaron así:
\begin{itemize}
  \item \textit{Suma de bytes.} 9Dh + 83h = 120h. C = 1 porque, sin signo, 157 + 131 = 288 supera 255. O = 1 porque, con signo, ($-99$) + ($-125$) = $-224$ no cabe en $-128$ a $127$ (dos negativos dan un resultado positivo). A = 1 porque Dh + 3h = 10h acarrea al nibble alto. S = 0 y Z = 0 porque 20h es positivo y distinto de cero, y P = 0 porque 20h tiene un solo uno.
  \item \textit{Suma de palabras.} 5938 + ($-5938$) = 0, es decir, 1732h + E8CEh = 10000h. Z = 1 y P = 1 porque el resultado de 16 bits es 0000h, y S = 0. C = 1 porque, sin signo, la suma rebasa FFFFh. O = 0 porque sumar dos números de signos opuestos nunca desborda. A = 1 porque 2h + Eh = 10h.
\end{itemize}
 
% ------------------------------------------------------------
\subsection{Segundo programa: resta con SUB y SBB}
 
\textbf{Objetivo.} Observar la resta de números de 1 y de 2 bytes, el comportamiento de las banderas C y A según el minuendo sea mayor o menor que el sustraendo, y la diferencia entre \texttt{SUB} y \texttt{SBB}.
 
\textbf{Funcionamiento.}
\begin{itemize}
  \item \texttt{CL = byte1 - byte2} (8 bits): 80 $-$ 69 = 11 (50h $-$ 45h = 0Bh).
  \item \texttt{DX = word1 - word2} (16 bits): 26437 $-$ 325 = 26112 (6745h $-$ 0145h = 6600h).
  \item El programa no copia los resultados a memoria: quedan en los registros CL y DX.
\end{itemize}
 
\begin{figure}[H]
    \centering
    \includegraphics[width=0.86\textwidth]{img/15.png}
    \caption{Código del segundo programa en el editor de EMU8086.}
    \label{fig: primerP_modificado.asm}
\end{figure}
 
\begin{table}[htbp]
  \centering
  \caption{Datos del segundo programa en decimal y en hexadecimal.}
  \label{tab:p2-datos}
  \begin{tabular}{|c|c|c|c|}
    \hline
    \textbf{Constante} & \textbf{Decimal} & \textbf{Hexadecimal} & \textbf{Papel} \\
    \hline
    \texttt{byte1} & 80 & 50h & minuendo (8 bits) \\
    \hline
    \texttt{byte2} & 69 & 45h & sustraendo (8 bits) \\
    \hline
    \texttt{word1} & 26437 & 6745h & minuendo (16 bits) \\
    \hline
    \texttt{word2} & 325 & 0145h & sustraendo (16 bits) \\
    \hline
  \end{tabular}
\end{table}
 
\textbf{Procedimiento.}
\begin{itemize}
  \item Se digitó y emuló el programa, anotando los minuendos y sustraendos en hexadecimal (Tabla~\ref{tab:p2-datos}).
  \item Se ejecutó paso a paso observando los resultados y las banderas (\textbf{1.ª prueba}, \texttt{SUB}).
  \item Se invirtió el orden de la operación, de modo que el minuendo pasa a ser sustraendo y viceversa (\textbf{2.ª prueba}, \texttt{SUB}):
\end{itemize}
 
\begin{figure}[H]
    \centering
    \includegraphics[width=0.86\textwidth]{img/15(1).png}
    \caption{Orden de operación invertido en el segundo programa.}
    \label{fig: segundoP_invertido.asm}
\end{figure}
 
\begin{itemize}
  \item Se restauró el programa original y se sustituyeron las instrucciones \texttt{SUB} por \texttt{SBB}, repitiendo las dos pruebas (\textbf{3.ª y 4.ª prueba}):
\end{itemize}
 
\begin{figure}[H]
    \centering
    \includegraphics[width=0.86\textwidth]{img/15(2).png}
    \caption{Sustitución de las instrucciones \textit{SUB} por \textit{SBB}.}
    \label{fig: segundoP_invertido.asm}
\end{figure}
 
\begin{table}[H]
  \centering
  \small
  \setlength{\tabcolsep}{3pt}
  \caption{Tabla 2: resultados de las restas con \texttt{SUB} y \texttt{SBB} (C = 0 al inicio).}
  \label{tab:p2-tabla2}
  \begin{tabular}{|c|l|c|c|c|c|c|c|c|c|c|}
    \hline
    \textbf{Prueba} & \textbf{Instrucción} & \textbf{Minuendo} & \textbf{Sustraendo} & \textbf{Diferencia} & \textbf{O} & \textbf{S} & \textbf{Z} & \textbf{A} & \textbf{P} & \textbf{C} \\
    \hline
    1.ª \texttt{SUB} & \texttt{sub cl, byte2} & 50h & 45h & 0Bh & 0 & 0 & 0 & 1 & 0 & 0 \\
    \hline
    1.ª \texttt{SUB} & \texttt{sub dx, word2} & 6745h & 0145h & 6600h & 0 & 0 & 0 & 0 & 1 & 0 \\
    \hline
    2.ª \texttt{SUB} & \texttt{sub cl, byte1} & 45h & 50h & F5h & 0 & 1 & 0 & 0 & 1 & 1 \\
    \hline
    2.ª \texttt{SUB} & \texttt{sub dx, word1} & 0145h & 6745h & 9A00h & 0 & 1 & 0 & 0 & 1 & 1 \\
    \hline
    3.ª \texttt{SBB} & \texttt{sbb cl, byte2} & 50h & 45h & 0Bh & 0 & 0 & 0 & 1 & 0 & 0 \\
    \hline
    3.ª \texttt{SBB} & \texttt{sbb dx, word2} & 6745h & 0145h & 6600h & 0 & 0 & 0 & 0 & 1 & 0 \\
    \hline
    4.ª \texttt{SBB} & \texttt{sbb cl, byte1} & 45h & 50h & F5h & 0 & 1 & 0 & 0 & 1 & 1 \\
    \hline
    4.ª \texttt{SBB} & \texttt{sbb dx, word1} & 0145h & 6745h & 99FFh & 0 & 1 & 0 & 1 & 1 & 1 \\
    \hline
  \end{tabular}
\end{table}


\subsubsection{Acarreo activado manualmente}
 
Se repitieron las pruebas activando manualmente la bandera C antes de cada resta (equivale a insertar \texttt{stc} antes de cada \texttt{sub} o \texttt{sbb}).
 
\begin{table}[htbp]
  \centering
  \small
  \setlength{\tabcolsep}{3pt}
  \caption{Restas con C = 1 antes de cada operación.}
  \label{tab:p2-acarreo}
  \begin{tabular}{|l|c|c|c|c|c|c|c|c|c|}
    \hline
    \textbf{Instrucción} & \textbf{Minuendo} & \textbf{Sustraendo} & \textbf{Diferencia} & \textbf{O} & \textbf{S} & \textbf{Z} & \textbf{A} & \textbf{P} & \textbf{C} \\
    \hline
    \texttt{sub cl, byte2} & 50h & 45h & 0Bh & 0 & 0 & 0 & 1 & 0 & 0 \\
    \hline
    \texttt{sub dx, word2} & 6745h & 0145h & 6600h & 0 & 0 & 0 & 0 & 1 & 0 \\
    \hline
    \texttt{sub cl, byte1} & 45h & 50h & F5h & 0 & 1 & 0 & 0 & 1 & 1 \\
    \hline
    \texttt{sub dx, word1} & 0145h & 6745h & 9A00h & 0 & 1 & 0 & 0 & 1 & 1 \\
    \hline
    \texttt{sbb cl, byte2} & 50h & 45h & 0Ah & 0 & 0 & 0 & 1 & 1 & 0 \\
    \hline
    \texttt{sbb dx, word2} & 6745h & 0145h & 65FFh & 0 & 0 & 0 & 1 & 1 & 0 \\
    \hline
    \texttt{sbb cl, byte1} & 45h & 50h & F4h & 0 & 1 & 0 & 0 & 0 & 1 \\
    \hline
    \texttt{sbb dx, word1} & 0145h & 6745h & 99FFh & 0 & 1 & 0 & 1 & 1 & 1 \\
    \hline
  \end{tabular}
\end{table}

\begin{figure}[H]
    \centering
    \includegraphics[width=0.86\textwidth]{img/15(3).png}
    \caption{Registros y banderas después de las restas con \texttt{SUB}.}
    \label{fig: segundoP.asm}
\end{figure}

\begin{figure}[H]
    \centering
    \includegraphics[width=0.86\textwidth]{img/15(4).png}
    \caption{Registros y banderas con \texttt{SUB} en orden invertido.}
    \label{fig: segundoP_invertido.asm}
\end{figure}

\begin{figure}[H]
    \centering
    \includegraphics[width=0.86\textwidth]{img/15(5).png}
    \caption{Registros y banderas con \texttt{SBB} en el orden original.}
    \label{fig: segundoP_sbb.asm}
\end{figure}
 
\begin{figure}[H]
    \centering
    \includegraphics[width=0.86\textwidth]{img/15(6).png}
    \caption{Registros y banderas con \texttt{SBB} en orden invertido.}
    \label{fig: segundoP_sbb_invertido.asm}
\end{figure}
 
\textbf{Observaciones.}
\begin{itemize}
  \item Con \texttt{SUB} el acarreo previo \textbf{no cambia} ninguna operación (Tabla~\ref{tab:p2-acarreo} es igual a la Tabla~\ref{tab:p2-tabla2} en las filas \texttt{SUB}).
  \item Con \texttt{SBB} sí cambia: cada diferencia disminuye en 1 (0Bh pasa a 0Ah, 6600h pasa a 65FFh, F5h pasa a F4h), porque la instrucción resta además el préstamo.
  \item Con \texttt{SBB} en orden invertido y C = 0 al inicio, la segunda resta ya da 99FFh en lugar de 9A00h: la primera resta dejó C = 1 (préstamo) y la segunda lo recibió.
\end{itemize}
 
% ------------------------------------------------------------
\subsection{Tercer programa: multiplicación con MUL e IMUL}
 
\textbf{Objetivo.} Observar la multiplicación de dos números de 1 byte, en qué registros se guardan los datos y el producto, y la diferencia entre multiplicar sin signo (\texttt{MUL}) y con signo (\texttt{IMUL}).
 
\textbf{Funcionamiento.}
\begin{itemize}
  \item AH = AAh y AL = FFh. El valor de AH se pierde, porque el producto de 16 bits sobrescribe todo AX.
  \item CL = 06h (multiplicador).
  \item \texttt{mul cl}: AX = AL $\times$ CL = FFh $\times$ 06h = 05FAh (255 $\times$ 6 = 1530).
\end{itemize}
 
\begin{figure}[H]
    \centering
    \includegraphics[width=0.86\textwidth]{img/16.png}
    \caption{Código del tercer programa en el editor de EMU8086}
    \label{fig: tercerP.asm}
\end{figure}
 
\textbf{Respuestas a las preguntas de la guía.}
\begin{itemize}
  \item \textit{¿Qué registros están involucrados?} AL (multiplicando, implícito), CL (multiplicador) y AX (producto). AH solo se carga con AAh, pero no participa.
  \item \textit{¿Qué datos se operan?} FFh y 06h.
  \item \textit{¿Dónde se almacena el resultado?} En AX (AH = parte alta, AL = parte baja).
\end{itemize}
 
\textbf{Procedimiento.} Se ejecutó paso a paso con el dato 06 y con \texttt{MUL} (Figura~\ref{fig:p3-mul}); luego con CL = $-06$ (Figura~\ref{fig:p3-mul-neg}); después se cambió \texttt{mul cl} por \texttt{imul cl} y se probó con 06 (Figura~\ref{fig:p3-imul}) y con $-06$ (Figura~\ref{fig:p3-imul-neg}).
 
\begin{table}[htbp]
  \centering
  \small
  \setlength{\tabcolsep}{3pt}
  \caption{Tabla 3: multiplicación con \texttt{MUL} e \texttt{IMUL}.}
  \label{tab:p3-tabla3}
  \begin{tabular}{|l|c|c|c|c|c|}
    \hline
    \textbf{Instrucción} & \textbf{Dato} & \textbf{Multiplicando (AL)} & \textbf{Multiplicador (CL)} & \textbf{Producto (AX)} & \textbf{C/O} \\
    \hline
    \texttt{MUL} & 06 & FFh = 255 & 06h = 6 & 05FAh = 1530 & 1/1 \\
    \hline
    \texttt{MUL} & $-06$ & FFh = 255 & FAh = 250 & F906h = 63750 & 1/1 \\
    \hline
    \texttt{IMUL} & 06 & FFh = $-1$ & 06h = 6 & FFFAh = $-6$ & 0/0 \\
    \hline
    \texttt{IMUL} & $-06$ & FFh = $-1$ & FAh = $-6$ & 0006h = 6 & 0/0 \\
    \hline
  \end{tabular}
 
  \smallskip
  {\small La columna C/O no la pide la guía, pero ayuda en el análisis. Las demás banderas quedan indefinidas.}
\end{table}

\begin{figure}[H]
    \centering
    \includegraphics[width=0.86\textwidth]{img/16(1).png}
    \caption{Resultado de \texttt{mul cl} con CL = 06.}
    \label{fig: tercerP.asm}
\end{figure}
 
\begin{figure}[H]
    \centering
    \includegraphics[width=0.86\textwidth]{img/16(2).png}
    \caption{Resultado de \texttt{mul cl} con CL = $-06$ (FAh).}
    \label{fig: tercerP_negative.asm}
\end{figure}
 
\begin{figure}[H]
    \centering
    \includegraphics[width=0.86\textwidth]{img/16(3).png}
    \caption{Resultado de \texttt{imul cl} con CL = 06.}
    \label{fig: tercerP_imul.asm}
\end{figure}
 
\begin{figure}[H]
    \centering
    \includegraphics[width=0.86\textwidth]{img/16(4).png}
    \caption{Resultado de \texttt{imul cl} con CL = $-06$ (FAh).}
    \label{fig: tercerP_imul_negative.asm}
\end{figure}
 
\textbf{Observaciones.} Los datos son los mismos bits en las cuatro pruebas (FFh y 06h o FAh), pero \texttt{MUL} los interpreta como números sin signo e \texttt{IMUL} como números con signo, y por eso los productos son distintos. Con \texttt{MUL} la parte alta del producto (AH $\neq$ 0) activa C y O. Con \texttt{IMUL} el producto con signo cabe en AL, así que C y O quedan en 0.
 
% ------------------------------------------------------------
\subsection{Cuarto programa: división con DIV e IDIV}
 
\textbf{Objetivo.} Observar la división de un número de 16 bits entre uno de 8 bits, en qué registros quedan el cociente y el residuo, y la diferencia entre dividir sin signo (\texttt{DIV}) y con signo (\texttt{IDIV}).
 
\textbf{Funcionamiento.}
\begin{itemize}
  \item AX = 123 = 007Bh (dividendo de 16 bits).
  \item CL = 05h (divisor, el dato).
  \item \texttt{div cl}: AL = cociente = 123 $\div$ 5 = 24 (18h) y AH = residuo = 3 (03h). AX queda en 0318h.
\end{itemize}
 
\begin{figure}[H]
    \centering
    \includegraphics[width=0.86\textwidth]{img/17.png}
    \caption{Código del cuarto programa en el editor de EMU8086.}
    \label{fig: cuartoP.asm}
\end{figure}
  
\textbf{Respuestas a las preguntas de la guía.}
\begin{itemize}
  \item \textit{¿Qué registros están involucrados?} AX (dividendo), CL (divisor), AL (cociente) y AH (residuo).
  \item \textit{¿Qué datos se operan?} 123 (007Bh) y 05h.
  \item \textit{¿Dónde se almacena el resultado?} En AX: el cociente en AL y el residuo en AH.
\end{itemize}
 
\textbf{Procedimiento.} Se siguió el mismo procedimiento del tercer programa: \texttt{DIV} con 05 y con $-05$ (Figuras~\ref{fig:p4-div} y~\ref{fig:p4-div-neg}); luego se cambió a \texttt{IDIV} y se probó con 05 y con $-05$ (Figuras~\ref{fig:p4-idiv} y~\ref{fig:p4-idiv-neg}).
 
\begin{table}[htbp]
  \centering
  \small
  \setlength{\tabcolsep}{3pt}
  \caption{Tabla 4: división con \texttt{DIV} e \texttt{IDIV}.}
  \label{tab:p4-tabla4}
  \begin{tabular}{|l|c|c|c|c|c|}
    \hline
    \textbf{Instrucción} & \textbf{Dato} & \textbf{Dividendo (AX)} & \textbf{Divisor (CL)} & \textbf{Cociente (AL)} & \textbf{Residuo (AH)} \\
    \hline
    \texttt{DIV} & 05 & 007Bh = 123 & 05h = 5 & 18h = 24 & 03h = 3 \\
    \hline
    \texttt{DIV} & $-05$ & 007Bh = 123 & FBh = 251 & 00h = 0 & 7Bh = 123 \\
    \hline
    \texttt{IDIV} & 05 & 007Bh = $+123$ & 05h = $+5$ & 18h = $+24$ & 03h = $+3$ \\
    \hline
    \texttt{IDIV} & $-05$ & 007Bh = $+123$ & FBh = $-5$ & E8h = $-24$ & 03h = $+3$ \\
    \hline
  \end{tabular}
 
  \smallskip
  {\small Las banderas quedan indefinidas después de \texttt{DIV} e \texttt{IDIV}.}
\end{table}
 
\begin{itemize}
    \item \textit{No se adjuntan imágenes de los resultados de las modificaciones del cuarto programa porque el comportamiento es muy similar al del tercer programa.} 
\end{itemize}

\subsubsection{Divisor igual a cero}
 
Se cambió el dato por 00h (\texttt{mov cl, 00h}) y se ejecutó el programa.
 
\begin{figure}[H]
    \centering
    \includegraphics[width=0.91\textwidth]{img/18.png}
    \caption{Error del emulador al dividir entre 00h.}
    \label{fig: cuartoP_divisorCero.asm}
\end{figure}
 
\textbf{Observaciones.} Al ejecutar \texttt{div cl} con CL = 00h el programa \textbf{no termina normalmente}: el emulador muestra un error de división (\textit{divide error}) y se detiene. La división entre cero no está definida, y el procesador responde generando la interrupción de tipo 0 en lugar de entregar un cociente.
 
% ============================================================
\section{Análisis de resultados}
 
\begin{enumerate}
  \renewcommand{\labelenumi}{\alph{enumi})}
 
  \item \textbf{Primer programa.}
  \begin{itemize}
    \item \textbf{Banderas con los datos originales.}
    \begin{itemize}
      \item \textit{Suma de bytes} (7Dh + 96h = 113h, AL = 13h). C = 1 porque 125 + 150 = 275 supera 255. A = 1 porque Dh + 6h = 13h acarrea del bit 3 al 4. O = 0 porque 7Dh es positivo ($+125$) y 96h es negativo ($-106$), y sumar números de signo distinto no desborda ($125 - 106 = 19$). S = 0 (el bit 7 de 13h es 0), Z = 0 y P = 0 (13h = 0001\,0011b tiene tres unos).
      \item \textit{Suma de palabras} (20A3h + 5938h = 79DBh). C = 0 porque no se pasa de FFFFh. A = 0 porque 3h + 8h = Bh no acarrea. O = 0 porque dos positivos dan 79DBh, que sigue siendo positivo (menor que 8000h). S = 0, Z = 0 y P = 1 (el byte bajo DBh = 1101\,1011b tiene seis unos).
    \end{itemize}
    \item \textbf{Banderas con los datos del paso 12.} 9Dh + 83h = 120h, con C = 1 (157 + 131 = 288 > 255), O = 1 ($-99 - 125 = -224$ no cabe en un byte con signo) y A = 1. En 16 bits, 1732h + E8CEh = 10000h, con Z = 1 y P = 1 (el resultado de 16 bits es 0000h), C = 1 y O = 0 (sumar un número y su opuesto no desborda).
    \item \textbf{¿Por qué cambiaron los resultados con \texttt{ADC} en lugar de \texttt{ADD}?} \texttt{ADD} calcula $\text{destino} + \text{fuente}$ y \texttt{ADC} calcula $\text{destino} + \text{fuente} + C$. Con C = 1 al inicio: 7Dh + 96h + 1 = 114h, es decir AL = 14h (en lugar de 13h). Como la primera suma vuelve a dejar C = 1, la segunda es 20A3h + 5938h + 1 = 79DCh (en lugar de 79DBh). Por eso \texttt{ADC} se usa para encadenar sumas de varios bytes: arrastra el acarreo de la parte baja a la parte alta.
  \end{itemize}
 
  \item \textbf{Segundo programa.}
  \begin{itemize}
    \item \textbf{C y A según el minuendo sea mayor o menor.} En la resta, C = 1 indica que hubo \emph{préstamo}: el minuendo es menor que el sustraendo y se tuvo que <<pedir prestado>> un 1 al bit siguiente al más significativo. Con 50h $-$ 45h (minuendo mayor) no hay préstamo y C = 0. Con 45h $-$ 50h (minuendo menor) se calcula (100h + 45h) $-$ 50h = F5h y C = 1. De forma análoga A = 1 cuando el nibble bajo del minuendo es menor que el del sustraendo: en 50h $-$ 45h, 0h < 5h, se pide prestado al nibble alto y A = 1; en 45h $-$ 50h, 5h $\geq$ 0h y A = 0 (el préstamo ocurre en el nibble alto, 4h < 5h, y por eso C = 1). Las restas de 16 bits siguen la misma lógica: 6745h $-$ 0145h no pide prestado (C = 0), mientras que 0145h $-$ 6745h = 10145h $-$ 6745h = 9A00h sí lo pide (C = 1).
    \item \textbf{¿Por qué cambiaron los resultados con \texttt{SBB} en lugar de \texttt{SUB}?} \texttt{SBB} calcula $\text{destino} - \text{fuente} - C$, es decir, resta además el préstamo anterior. Con C = 0 la primera resta es igual que con \texttt{SUB} (0Bh y 6600h). En el orden invertido, la resta de bytes deja C = 1 y la siguiente es 0145h $-$ 6745h $-$ 1 = 99FFh (en lugar de 9A00h). Si C se activa a mano antes de cada resta, todos los resultados disminuyen en 1: 50h $-$ 45h $-$ 1 = 0Ah, 6745h $-$ 0145h $-$ 1 = 65FFh y 45h $-$ 50h $-$ 1 = F4h. Así, \texttt{SBB} permite encadenar restas de varios bytes.
  \end{itemize}
 
  \item \textbf{Tercer programa.}
  \begin{itemize}
    \item \textbf{¿Cómo se especifica que se trabaja con números de un byte?} Por el tamaño del operando que acompaña a la instrucción: \texttt{mul cl} tiene un operando de 8 bits (CL) y, por tanto, el otro factor está implícito en AL. Si el operando fuera de 16 bits (por ejemplo \texttt{mul cx}), se multiplicaría AX $\times$ CX.
    \item \textbf{Registros que se usan.} El multiplicando va en AL, el multiplicador en un registro o posición de memoria de 8 bits (CL) y el producto de 16 bits queda en AX. (En multiplicaciones de 16 bits el producto queda en el par DX:AX.)
    \item \textbf{¿Por qué cambiaron los resultados con \texttt{IMUL} en lugar de \texttt{MUL}?} Los bits son los mismos, pero la interpretación es distinta. Con 06: \texttt{MUL} calcula 255 $\times$ 6 = 1530 = 05FAh e \texttt{IMUL} calcula $(-1) \times 6 = -6$ = FFFAh. Con $-06$ (FAh): \texttt{MUL} calcula 255 $\times$ 250 = 63750 = F906h e \texttt{IMUL} calcula $(-1) \times (-6) = 6$ = 0006h. Además, con \texttt{MUL} se activan C y O porque AH $\neq$ 0, mientras que con \texttt{IMUL} se quedan en 0 porque el producto con signo cabe en AL.
  \end{itemize}
 
  \item \textbf{Cuarto programa.}
  \begin{itemize}
    \item \textbf{¿Cómo se especifica que se trabaja con números de un byte?} También por el tamaño del divisor: \texttt{div cl} tiene un divisor de 8 bits, por lo que el dividendo es AX. Con un divisor de 16 bits (por ejemplo \texttt{div cx}) el dividendo sería DX:AX.
    \item \textbf{Registros que se usan.} Dividendo de 16 bits en AX, divisor de 8 bits en CL, cociente en AL y residuo en AH. (Con divisor de 16 bits: dividendo en DX:AX, cociente en AX y residuo en DX.)
    \item \textbf{¿Por qué cambiaron los resultados con \texttt{IDIV} en lugar de \texttt{DIV}?} Con 05 ambas instrucciones dan lo mismo, porque todos los datos son positivos: 123 $\div$ 5 = 24 con residuo 3 (0318h). Con $-05$ cambian. \texttt{DIV} ve CL = FBh como 251 y calcula 123 $\div$ 251 = 0 con residuo 123 (AX = 7B00h). \texttt{IDIV} ve CL como $-5$ y calcula 123 $\div$ ($-5$) = $-24$ con residuo $+3$ (AL = E8h, AH = 03h, AX = 03E8h); se comprueba con $(-24)(-5) + 3 = 123$. El residuo de \texttt{IDIV} lleva el signo del dividendo.
    \item \textbf{¿Por qué la última modificación generó un error?} Porque con CL = 00h se intenta dividir entre cero, y esa operación no está definida. El procesador genera la excepción de división (interrupción tipo 0) y el emulador detiene el programa. El mismo error ocurre si el cociente no cabe en AL (más de FFh con \texttt{DIV}, o fuera de $-128$ a $127$ con \texttt{IDIV}).
  \end{itemize}
\end{enumerate}
 
% ============================================================
\section{Investigación complementaria}
 
\subsection{Directivas EQU y DB}
 
\begin{table}[htbp]
  \centering
  \caption{Comparación entre las directivas \texttt{EQU} y \texttt{DB}.}
  \label{tab:equ-db}
  \begin{tabular}{|p{3.4cm}|p{5.2cm}|p{5.2cm}|}
    \hline
    \textbf{Característica} & \textbf{\texttt{EQU}} & \textbf{\texttt{DB}} \\
    \hline
    Qué hace & Asigna un nombre a un valor constante. & Reserva y asigna valor inicial a bytes de memoria. \\
    \hline
    Memoria & No reserva memoria: el ensamblador reemplaza el nombre por el valor. & Reserva 1 byte por cada valor. \\
    \hline
    Uso en la instrucción & Operando inmediato (\texttt{mov al, byte1} se ensambla como \texttt{mov al, 7Dh}). & Operando de memoria (direccionamiento directo): \texttt{mov al, dato}. \\
    \hline
    Modificable al ejecutar & No, es una constante. & Sí, se puede leer y escribir (\texttt{mov dato, al}). \\
    \hline
    Dirección & No tiene. & La etiqueta equivale al desplazamiento donde queda el dato. \\
    \hline
  \end{tabular}
\end{table}
 
\subsection{Multiplicación de dos números de dos bytes}
 
Se usa la forma de 16 bits de \texttt{MUL} (o \texttt{IMUL} si hay signo): el multiplicando va en AX, el multiplicador en un registro de 16 bits y el producto de 32 bits queda en el par DX:AX (DX = parte alta, AX = parte baja).
 
\begin{minted}[linenos,frame=lines,bgcolor=bg,breaklines,fontsize=\small]{nasm}
    mov ax, 1234h       ; multiplicando
    mov bx, 0100h       ; multiplicador
    mul bx              ; DX:AX = 1234h x 0100h = 0012 3400h
\end{minted}
 
Como DX $\neq$ 0, las banderas C y O quedan en 1. También puede hacerse con registros de 8 bits mediante productos parciales: $(a_H 2^8 + a_L)(b_H 2^8 + b_L) = a_H b_H 2^{16} + (a_H b_L + a_L b_H) 2^8 + a_L b_L$, con cuatro \texttt{MUL} de 8 bits y sumas encadenadas con \texttt{ADC}.
 
\subsection{División de dos números de dos bytes}
 
Se usa la forma de 16 bits de \texttt{DIV} (o \texttt{IDIV}): el dividendo de 32 bits va en DX:AX, el divisor en un registro de 16 bits, el cociente queda en AX y el residuo en DX. Si el dividendo solo tiene 16 bits, hay que completar DX: con 0 para \texttt{DIV} (\texttt{xor dx, dx}) o con la extensión de signo para \texttt{IDIV} (\texttt{cwd}).
 
\begin{minted}[linenos,frame=lines,bgcolor=bg,breaklines,fontsize=\small]{nasm}
    mov ax, 5000        ; dividendo (16 bits)
    xor dx, dx          ; DX = 0, completa el dividendo a 32 bits
    mov bx, 7           ; divisor
    div bx              ; AX = 02CAh (714), DX = 0002h (residuo)
\end{minted}
 
Se verifica con $714 \times 7 + 2 = 5000$. Igual que en 8 bits, hay error de división si el divisor es 0 o si el cociente no cabe en AX.
 
\subsection{Suma de dos números de 4 bytes}
 
Se suman A = 0012FFFFh y B = 00340001h. Como ningún registro del 8086 guarda 32 bits, la suma se hace por partes y el acarreo de cada parte se arrastra a la siguiente con \texttt{ADC}: la parte menos significativa se suma con \texttt{ADD} y las demás con \texttt{ADC}. En ambos casos el resultado esperado es R = 00470000h.
 
\subsubsection{Con registros de 16 bits}
 
Se suman dos palabras: primero la baja con \texttt{ADD} y luego la alta con \texttt{ADC}. Entre ambas solo hay instrucciones \texttt{MOV}, que no alteran C.
 
\begin{figure}[htbp]
  \centering
  \begin{tikzpicture}[
      term/.style={draw, rounded corners=9pt, minimum width=3.2cm, minimum height=0.65cm, align=center, font=\footnotesize},
      caja/.style={draw, minimum width=7.6cm, minimum height=0.65cm, align=center, font=\footnotesize},
      flecha/.style={->, >=stealth}]
    \node[term] (n0) at (0,0)     {INICIO};
    \node[caja] (n1) at (0,-1.1)  {AX $\leftarrow$ A$_{\mathrm{bajo}}$ (palabra baja de A)};
    \node[caja] (n2) at (0,-2.2)  {AX $\leftarrow$ AX $+$ B$_{\mathrm{bajo}}$ \quad (\texttt{ADD}, genera C)};
    \node[caja] (n3) at (0,-3.3)  {R$_{\mathrm{bajo}}$ $\leftarrow$ AX};
    \node[caja] (n4) at (0,-4.4)  {AX $\leftarrow$ A$_{\mathrm{alto}}$ (palabra alta de A)};
    \node[caja] (n5) at (0,-5.5)  {AX $\leftarrow$ AX $+$ B$_{\mathrm{alto}}$ $+$ C \quad (\texttt{ADC})};
    \node[caja] (n6) at (0,-6.6)  {R$_{\mathrm{alto}}$ $\leftarrow$ AX};
    \node[term] (n7) at (0,-7.7)  {FIN};
    \foreach \a/\b in {n0/n1,n1/n2,n2/n3,n3/n4,n4/n5,n5/n6,n6/n7}
      \draw[flecha] (\a) -- (\b);
  \end{tikzpicture}
  \caption{Diagrama de flujo de la suma de dos números de 4 bytes con registros de 16 bits.}
  \label{fig:flujo16}
\end{figure}
 
\begin{figure}[H]
    \centering
    \includegraphics[width=0.91\textwidth]{img/19.png}
    \caption{Ejecución de la suma de 4 bytes con registros de 16 bits.}
    \label{fig: suma4bytes.asm}
\end{figure}
  
\subsubsection{Con registros de 8 bits}
 
Se suman cuatro bytes, del menos al más significativo: el primero con \texttt{ADD} y los otros tres con \texttt{ADC}.
 
\begin{figure}[htbp]
  \centering
  \begin{tikzpicture}[
      term/.style={draw, rounded corners=9pt, minimum width=3.2cm, minimum height=0.55cm, align=center, font=\footnotesize},
      caja/.style={draw, minimum width=7.4cm, minimum height=0.55cm, align=center, font=\footnotesize},
      flecha/.style={->, >=stealth}]
    \node[term] (n0)  at (0,0)      {INICIO};
    \node[caja] (n1)  at (0,-0.95)  {AL $\leftarrow$ A$_0$ (byte menos significativo de A)};
    \node[caja] (n2)  at (0,-1.9)   {AL $\leftarrow$ AL $+$ B$_0$ \quad (\texttt{ADD}, genera C)};
    \node[caja] (n3)  at (0,-2.85)  {R$_0$ $\leftarrow$ AL};
    \node[caja] (n4)  at (0,-3.8)   {AL $\leftarrow$ A$_1$};
    \node[caja] (n5)  at (0,-4.75)  {AL $\leftarrow$ AL $+$ B$_1$ $+$ C \quad (\texttt{ADC})};
    \node[caja] (n6)  at (0,-5.7)   {R$_1$ $\leftarrow$ AL};
    \node[caja] (n7)  at (0,-6.65)  {AL $\leftarrow$ A$_2$};
    \node[caja] (n8)  at (0,-7.6)   {AL $\leftarrow$ AL $+$ B$_2$ $+$ C \quad (\texttt{ADC})};
    \node[caja] (n9)  at (0,-8.55)  {R$_2$ $\leftarrow$ AL};
    \node[caja] (n10) at (0,-9.5)   {AL $\leftarrow$ A$_3$ (byte más significativo de A)};
    \node[caja] (n11) at (0,-10.45) {AL $\leftarrow$ AL $+$ B$_3$ $+$ C \quad (\texttt{ADC})};
    \node[caja] (n12) at (0,-11.4)  {R$_3$ $\leftarrow$ AL};
    \node[term] (n13) at (0,-12.35) {FIN};
    \foreach \a/\b in {n0/n1,n1/n2,n2/n3,n3/n4,n4/n5,n5/n6,n6/n7,n7/n8,n8/n9,n9/n10,n10/n11,n11/n12,n12/n13}
      \draw[flecha] (\a) -- (\b);
  \end{tikzpicture}
  \caption{Diagrama de flujo de la suma de dos números de 4 bytes con registros de 8 bits.}
  \label{fig:flujo8}
\end{figure}

\begin{figure}[H]
    \centering
    \includegraphics[width=0.91\textwidth]{img/19(1).png}
    \caption{Ejecución de la suma de 4 bytes con registros de 8 bits.}
    \label{fig: suma4bytes_8bits.asm}
\end{figure}
 
% ============================================================
\section{Conclusiones}
 
\begin{itemize}
  \item \texttt{ADD} y \texttt{SUB} solo usan los operandos, mientras que \texttt{ADC} y \texttt{SBB} incluyen además el acarreo o préstamo previo. Por eso estas últimas permiten operar números más grandes que un registro, encadenando partes de menor a mayor peso.
  \item La bandera C señala el desbordamiento de números \emph{sin signo} (acarreo en la suma, préstamo en la resta) y la bandera O el de números \emph{con signo}. Un mismo resultado puede activar una, la otra, ambas o ninguna: 9Dh $+$ 83h activó las dos, mientras que 7Dh $+$ 96h activó solo C.
  \item En la resta, C = 1 indica préstamo (minuendo menor que sustraendo) y A = 1 indica préstamo entre nibbles; el resultado queda en complemento a 2 (45h $-$ 50h = F5h).
  \item El tamaño de una multiplicación o división lo fija el operando: con un operando de 8 bits el otro dato está implícito en AL (\texttt{MUL}) o en AX (\texttt{DIV}); con uno de 16 bits se usa AX o DX:AX. El producto de \texttt{MUL} ocupa el doble de bits de los factores.
  \item \texttt{MUL} y \texttt{DIV} interpretan los datos sin signo, e \texttt{IMUL} e \texttt{IDIV} con signo. Con los mismos bits (FFh y FAh) los resultados difieren: 1530 frente a $-6$ en la multiplicación, y 0 frente a $-24$ en la división. En \texttt{IDIV} el residuo lleva el signo del dividendo.
  \item Dividir entre cero, o con un cociente que no cabe en el registro de destino, genera un error de ejecución (interrupción tipo 0).
  \item \texttt{EQU} define constantes que el ensamblador sustituye en el código, mientras que \texttt{DB} reserva memoria con una dirección y un valor que se puede modificar.
  \item Las instrucciones aritméticas modifican las banderas (O, S, Z, A, P, C) según el resultado, salvo \texttt{MUL}, \texttt{IMUL}, \texttt{DIV} e \texttt{IDIV}, que dejan indefinidas varias de ellas.
\end{itemize}
 
% ============================================================
\section{Bibliografía}
 
% Completa cada entrada con la edición, el año o la URL de la fuente que realmente consultaste.
\begin{itemize}
  \item Guía de laboratorio de Arquitectura de Computadoras, Práctica 5: Programación para el microprocesador INTEL, instrucciones aritméticas. Escuela Profesional de Ingeniería de Sistemas.
  \item Brey, B. B. \emph{Los microprocesadores Intel}. Pearson Educación.
  \item Intel Corporation. \emph{The 8086 Family User's Manual}.
  \item Documentación (ayuda) del emulador EMU8086.
\end{itemize}

\end{document}